\section{Discussion}
\label{sec:discussion}

\subsection{What the comparison establishes}

The benchmark supports a narrow systems claim.  On the fixed observations in this
development set, the offline pipeline occupies a substantially better
continuity--contamination operating point than the evaluated OC-SORT and BoT-SORT
configurations.  This is useful because the final product requires long identity
paths rather than a large collection of locally plausible fragments.  It is not a
claim of general MOT superiority: the observation set is domain-specific, the
evaluation is in-sample, and detector recall is outside the protocol.

The results also make the asymmetric objective concrete.  OC-SORT has the highest
pair precision, yet its 845 excess fragments make it operationally unsuitable.  The
production system reduces that count to 42, but the remaining nine contaminated
tracks are still serious errors.  Aggregate F1 is therefore a diagnostic rather
than an acceptance criterion; final trajectories remain subject to review.

\subsection{Residual failure mode}

The original failure audit separated two kinds of false link.  Geometrically
implausible joins resulted from bypassed foreground masking, endpoint-biased
appearance aggregation, and position-only use of a size-aware motion model.  The
revised evidence and conservative vetoes remove many of these cases.  The remaining
wrong joins cluster at crossings, visually near-identical sails, and crops containing
both competitors.  In those cases, time, position, scale, and color can all support
the wrong continuation.

This residual mode is not well addressed by another global scalar threshold.  An
abstention experiment removed each chosen edge in turn and compared the alternative
objective.  The resulting margin did not separate wrong from correct links: useful
thresholds rejected substantially more correct associations than errors.  A credible
next step would instead model the interaction explicitly---for example, by detecting
paired crossings and comparing pre- and post-overlap appearance jointly.  We leave
that extension unimplemented rather than tuning a development-set-specific margin.

\subsection{Scope and limitations}

The manually reconstructed videos are the development set, not a held-out test set.
The current tracker was frozen for the reported comparison, but the provenance of
the older numerical ILP tuning is incomplete.  Moreover, saved observations were
selected during manual reconstruction and their boxes were smoothed using the gold
trajectories.  The study consequently evaluates association conditional on a useful
set of boxes; it does not measure end-to-end detector performance and should not be
read as a MOTA or HOTA result.

The baseline configurations are also intentionally pragmatic.  BoT-SORT retains ECC
camera compensation but omits its pedestrian ReID model, whose learned representation
does not match windsurf sails.  A dedicated learned embedding might improve that
baseline, but training one would duplicate the domain adaptation that the compact
color descriptor is meant to avoid.  The pose model is documented as part of the
system and its geometry is visible in the produced crops, but it has no independent
cross-video quantitative benchmark.  Its image-level validation split may contain
nearby frames from the same source video.

Finally, the three-second candidate horizon encodes a practical trade-off.  It limits
false long-range joins and keeps the graph tractable, but a genuine disappearance
longer than that remains split.  The system therefore approaches its application
through conservative automation plus review, not through a guarantee of perfect
identity recovery.
